% AI-DLC Workflows 2.11 lifecycle (chapter 13): five phases, their stages, and the gates.
\begin{tikzpicture}[x=31.5mm, y=1mm,
    stg/.style={font=\sffamily\scriptsize, text width=29mm, align=left, anchor=north, inner sep=1pt, execute at begin node={\hyphenpenalty=10000}}]
  \foreach \i/\name/\tone/\n in {0/Initialization/notetone/3, 1/Ideation/todotone/7,
      2/Inception/decisiontone/9, 3/Construction/ruletone/7, 4/Operation/audittone/7} {
    \node[phase=\tone, text width=27mm] (p\i) at (\i,0) {\name\\[-1pt]{\normalfont\tiny\color{muted}\n\ stages}};
  }
  \foreach \i in {0,...,3} \pgfmathtruncatemacro{\j}{\i+1} \draw[flow] (p\i) -- (p\j);
  \node[stg] at (0,-7) {Workspace scaffold\\Workspace detection\\State initialization\\[2pt]\textit{no gates}};
  \node[stg] at (1,-7) {Intent capture\\Market research\\Feasibility\\Scope definition\\Team formation\\Rough mockups\\Approval \& handoff};
  \node[stg] at (2,-7) {Reverse engineering\\Practices discovery\\Requirements analysis\\User stories\\Refined mockups\\Domain design\\Units generation\\Contract design\\Delivery planning};
  \node[stg] at (3,-7) {Functional design\\NFR requirements\\NFR design\\Infrastructure design\\Code generation\\Build and test\\CI pipeline};
  \node[stg] at (4,-7) {Deployment pipeline\\Environment provisioning\\Deployment execution\\Observability setup\\Incident response\\Performance validation\\Feedback \& optimization};
  \node[font=\sffamily\scriptsize, text=muted, align=center] at (2,-52) {AI proposes a plan and asks $\rightarrow$ a person answers and approves at each gate $\rightarrow$ AI executes;\\every question, answer and approval is written to the intent's audit trail};
\end{tikzpicture}
